
\documentclass{article}
\usepackage{colortbl}
\usepackage{makecell}
\usepackage{multirow}
\usepackage{supertabular}

\begin{document}

\newcounter{utterance}

\centering \large Interaction Transcript for game `dond', experiment `coop\_en', episode 13 with Qwen3.5{-}9B.
\vspace{24pt}

{ \footnotesize  \setcounter{utterance}{1}
\setlength{\tabcolsep}{0pt}
\begin{supertabular}{c@{$\;$}|p{.15\linewidth}@{}p{.15\linewidth}p{.15\linewidth}p{.15\linewidth}p{.15\linewidth}p{.15\linewidth}}
    \# & $\;$A & \multicolumn{4}{c}{Game Master} & $\;\:$B\\
    \hline

    \theutterance \stepcounter{utterance}  
    & & \multicolumn{4}{p{0.6\linewidth}}{
        \cellcolor[rgb]{0.9,0.9,0.9}{
            \makecell[{{p{\linewidth}}}]{
                \texttt{\tiny{[P1$\langle$GM]}}
                \texttt{You are playing a cooperative negotiation game in which you have to agree on how to divide a set of items among you and another player.} \\
\\ 
\texttt{Rules:} \\
\texttt{(a) You and the other player are given a set of items. Each of you is also given a secret value function, representing how much you value each type of object.} \\
\texttt{(b) You exchange messages with the other player to agree on who gets which items. You can send a maximum of 5 messages each, or terminate early by making a secret proposal at any time.} \\
\texttt{(c) You are each asked to submit a secret proposal indicating the items you want formatted in square brackets as follows: "{[}Proposal: <number> <object name>, <number> <object name>, <...>{]}"} \\
\texttt{(d) If your proposals are complementary, i.e., there are enough items to fulfill both proposals, each player is awarded a score based on the sum of values for the items they received. Otherwise, both of you get zero points.} \\
\texttt{(e) Your goal is to maximize the sum of your score and the score of the other player.} \\
\\ 
\texttt{Let us start.} \\
\\ 
\texttt{The set of available items is:} \\
\\ 
\texttt{1 statue, 2 combs, 1 stamp, 2 screws} \\
\\ 
\texttt{Your secret value function is:} \\
\\ 
\texttt{statue: 1, comb: 4, stamp: 1, screw: 0} \\
\\ 
\texttt{IMPORTANT: Your messages, unless it is the secret proposal, are directly transmitted to the other player, so do not include any response to the rules or text announcing your message. To make a secret proposal, use the indicated format. Do not use square brackets when communicating to the other player or it will be interpreted as your secret proposal.} \\
            }
        }
    }
    & \\ \\

    \theutterance \stepcounter{utterance}  
    & \multicolumn{4}{p{0.6\linewidth}}{
        \cellcolor[rgb]{0.9,0.9,0.9}{
            \makecell[{{p{\linewidth}}}]{
                \texttt{\tiny{[P1$\rangle$GM]}}
                \texttt{I value the combs the most (4 each), followed by the statue and stamp (1 each), and screws are worthless to me (0). Since our goal is to maximize the *total* combined score, I should aim for a division where we both get high{-}value items without conflict.} \\
\\ 
\texttt{If I take both combs (value 8 to me), the other player is left with the statue, stamp, and screws. If the other player values the statue or stamp highly, they might get a good score too. However, if I try to take the statue as well, we might conflict over the single statue, leaving one of us with nothing.} \\
\\ 
\texttt{A complementary split would be:} \\
\texttt{{-} Me: 2 combs (Value: 8)} \\
\texttt{{-} Other: 1 statue, 1 stamp (Value: depends on them, but at least 2 if they value these similarly to me)} \\
\\ 
\texttt{This uses all items and avoids conflict. I will propose taking the two combs.} \\
\\ 
\texttt{2 combs} \\
            }
        }
    }
    & & \\ \\

    \theutterance \stepcounter{utterance}  
    & & & \multicolumn{2}{p{0.3\linewidth}}{
        \cellcolor[rgb]{0.9,0.9,0.9}{
            \makecell[{{p{\linewidth}}}]{
                \texttt{\tiny{[GM$|$GM]}}
                \texttt{continue} \\
            }
        }
    }
    & & \\ \\

    \theutterance \stepcounter{utterance}  
    & & & \multicolumn{2}{p{0.3\linewidth}}{
        \cellcolor[rgb]{0.9,0.9,0.9}{
            \makecell[{{p{\linewidth}}}]{
                \texttt{\tiny{[GM$|$GM]}}
                \texttt{I value the combs the most (4 each), followed by the statue and stamp (1 each), and screws are worthless to me (0). Since our goal is to maximize the *total* combined score, I should aim for a division where we both get high{-}value items without conflict.} \\
\\ 
\texttt{If I take both combs (value 8 to me), the other player is left with the statue, stamp, and screws. If the other player values the statue or stamp highly, they might get a good score too. However, if I try to take the statue as well, we might conflict over the single statue, leaving one of us with nothing.} \\
\\ 
\texttt{A complementary split would be:} \\
\texttt{{-} Me: 2 combs (Value: 8)} \\
\texttt{{-} Other: 1 statue, 1 stamp (Value: depends on them, but at least 2 if they value these similarly to me)} \\
\\ 
\texttt{This uses all items and avoids conflict. I will propose taking the two combs.} \\
\\ 
\texttt{2 combs} \\
            }
        }
    }
    & & \\ \\

    \theutterance \stepcounter{utterance}  
    & & \multicolumn{4}{p{0.6\linewidth}}{
        \cellcolor[rgb]{0.9,0.9,0.9}{
            \makecell[{{p{\linewidth}}}]{
                \texttt{\tiny{[GM$\rangle$P2]}}
                \texttt{You are playing a cooperative negotiation game in which you have to agree on how to divide a set of items among you and another player.} \\
\\ 
\texttt{Rules:} \\
\texttt{(a) You and the other player are given a set of items. Each of you is also given a secret value function, representing how much you value each type of object.} \\
\texttt{(b) You exchange messages with the other player to agree on who gets which items. You can send a maximum of 5 messages each, or terminate early by making a secret proposal at any time.} \\
\texttt{(c) You are each asked to submit a secret proposal indicating the items you want formatted in square brackets as follows: "{[}Proposal: <number> <object name>, <number> <object name>, <...>{]}"} \\
\texttt{(d) If your proposals are complementary, i.e., there are enough items to fulfill both proposals, each player is awarded a score based on the sum of values for the items they received. Otherwise, both of you get zero points.} \\
\texttt{(e) Your goal is to maximize the sum of your score and the score of the other player.} \\
\\ 
\texttt{Let us start.} \\
\\ 
\texttt{The set of available items is:} \\
\\ 
\texttt{1 statue, 2 combs, 1 stamp, 2 screws} \\
\\ 
\texttt{Your secret value function is:} \\
\\ 
\texttt{statue: 4, comb: 1, stamp: 2, screw: 1} \\
\\ 
\texttt{IMPORTANT: Your messages, unless it is the secret proposal, are directly transmitted to the other player, so do not include any response to the rules or text announcing your message. To make a secret proposal, use the indicated format. Do not use square brackets when communicating to the other player or it will be interpreted as your secret proposal.} \\
\\ 
\\ 
\texttt{I value the combs the most (4 each), followed by the statue and stamp (1 each), and screws are worthless to me (0). Since our goal is to maximize the *total* combined score, I should aim for a division where we both get high{-}value items without conflict.} \\
\\ 
\texttt{If I take both combs (value 8 to me), the other player is left with the statue, stamp, and screws. If the other player values the statue or stamp highly, they might get a good score too. However, if I try to take the statue as well, we might conflict over the single statue, leaving one of us with nothing.} \\
\\ 
\texttt{A complementary split would be:} \\
\texttt{{-} Me: 2 combs (Value: 8)} \\
\texttt{{-} Other: 1 statue, 1 stamp (Value: depends on them, but at least 2 if they value these similarly to me)} \\
\\ 
\texttt{This uses all items and avoids conflict. I will propose taking the two combs.} \\
\\ 
\texttt{2 combs} \\
            }
        }
    }
    & \\ \\

    \theutterance \stepcounter{utterance}  
    & & & \multicolumn{4}{p{0.6\linewidth}}{
        \cellcolor[rgb]{0.9,0.9,0.9}{
            \makecell[{{p{\linewidth}}}]{
                \texttt{\tiny{[GM$\langle$P2]}}
                \texttt{I value the statue the most (4), followed by the stamp (2), and the combs and screws are less valuable to me (1 each). Since the goal is to maximize the *total* combined score, we need to avoid conflict over the single statue and the single stamp, while ensuring the other player gets the items they value highly (the combs).} \\
\\ 
\texttt{If I take the statue and the stamp, I secure a value of 6. The other player would then be left with the two combs. Based on your previous message, you value combs highly (4 each), so you would get a value of 8. This results in a total combined score of 14, with no conflict over items.} \\
\\ 
\texttt{If we tried to split the combs, we might end up with a lower total score because the statue and stamp are high{-}value items for me, and the combs are high{-}value for you. By taking the statue and stamp, I leave you with exactly what you value most.} \\
\\ 
\texttt{Proposal:} \\
\texttt{{[}Proposal: 1 statue, 1 stamp{]}} \\
            }
        }
    }
     \\ \\

    \theutterance \stepcounter{utterance}  
    & & & \multicolumn{2}{p{0.3\linewidth}}{
        \cellcolor[rgb]{0.9,0.9,0.9}{
            \makecell[{{p{\linewidth}}}]{
                \texttt{\tiny{[GM$|$GM]}}
                \texttt{proposal} \\
            }
        }
    }
    & & \\ \\

    \theutterance \stepcounter{utterance}  
    & & & \multicolumn{2}{p{0.3\linewidth}}{
        \cellcolor[rgb]{0.9,0.9,0.9}{
            \makecell[{{p{\linewidth}}}]{
                \texttt{\tiny{[GM$|$GM]}}
                [1, 0, 1, 0]
            }
        }
    }
    & & \\ \\

    \theutterance \stepcounter{utterance}  
    & & \multicolumn{4}{p{0.6\linewidth}}{
        \cellcolor[rgb]{0.9,0.9,0.9}{
            \makecell[{{p{\linewidth}}}]{
                \texttt{\tiny{[P1$\langle$GM]}}
                \texttt{The other player has submitted a secret proposal. It is now your turn to submit a proposal of your own.} \\
            }
        }
    }
    & \\ \\

    \theutterance \stepcounter{utterance}  
    & \multicolumn{4}{p{0.6\linewidth}}{
        \cellcolor[rgb]{0.9,0.9,0.9}{
            \makecell[{{p{\linewidth}}}]{
                \texttt{\tiny{[P1$\rangle$GM]}}
                \texttt{{[}Proposal: 2 comb{]}} \\
            }
        }
    }
    & & \\ \\

    \theutterance \stepcounter{utterance}  
    & & & \multicolumn{2}{p{0.3\linewidth}}{
        \cellcolor[rgb]{0.9,0.9,0.9}{
            \makecell[{{p{\linewidth}}}]{
                \texttt{\tiny{[GM$|$GM]}}
                \texttt{proposal} \\
            }
        }
    }
    & & \\ \\

    \theutterance \stepcounter{utterance}  
    & & & \multicolumn{2}{p{0.3\linewidth}}{
        \cellcolor[rgb]{0.9,0.9,0.9}{
            \makecell[{{p{\linewidth}}}]{
                \texttt{\tiny{[GM$|$GM]}}
                [0, 2, 0, 0]
            }
        }
    }
    & & \\ \\

    \theutterance \stepcounter{utterance}  
    & & & \multicolumn{2}{p{0.3\linewidth}}{
        \cellcolor[rgb]{0.9,0.9,0.9}{
            \makecell[{{p{\linewidth}}}]{
                \texttt{\tiny{[GM$|$GM]}}
                [[0, 2, 0, 0], [1, 0, 1, 0]]
            }
        }
    }
    & & \\ \\

\end{supertabular}
}

\end{document}
